\chapter{Hygiene --- \gAasRsN{RS~II}}
\label{bkm:RefHeading32633609}Sebastian Gabriel [*]

\marginpar{\includegraphics[width=4cm]{images/AASdev20090228-img95.jpg}}

\minitoc



\gChangelog{
Changelog:

\begin{description}
    \item[2009-02-20] Fertig f\"ur Kurs 2009/02.
    \item[2009-02-25] Start Changelog.
    \item[2009-04-02] Übernahme aus Openoffice in \LaTeX
\end{description}
}

\section{Grundlagen}
\begin{multicols}{2}
\begin{description}
    \item[\gT{Infektion}] \index{Infektion}Eindringen von
    krankmachenden Keimen in den Organismus
    \item[\gT{Pathogenit\"at}]
    \index{Pathogenit\"at}F\"ahigkeit eines Mikroorganismus, eine Krankheit auszul\"osen
    \item[\gT{Resistenz}]
    \index{Resistenz}Widerstandsf\"ahigkeit eines Organismus gegen die  krankmachenden  Eigenschaften eines Erregers
    \item[\gT{Inkubationszeit}]
    \index{Inkubationszeit}symptomlose Zeit zwischen Infektion (Eintritt d. Erregers in den Organismus) und dem Auftreten
    erster Symptome
    \item[\gT{Epidemie}] \index{Epidemie}\"ortlich
    begrenztes, zeitlich vermehrtes Auftreten einer  Infektionserkrankung (Grippe, Salmonellen)
    \item[\gT{Pandemie}] \index{Pandemie}Epidemie die sich
    \"uber L\"ander u. Kontinente ausbreitet (Pest im Mittelalter)
    \item[\gZ{Endemie}] \index{Endemie}Auf kleinere Gebiete
    beschr\"anktes, zeitlich jedoch nicht begrenztes Auftreten einer Infektionserkrankung  (FSME,
    Malaria, Lepra)
    \item[\gT{Morbidit\"at}] \index{Morbidit\"at}Anteil der
    an XY Erkrankten einer bestimmten Bev\"olkerungsgruppe (z.B. 5 pro 100.000 EW)
    \item[\gT{Mortalit\"at}] \index{Mortalit\"at}Anteil der
    an XY Verstorbenen einer bestimmten Bev\"olkerungsgruppe (z.B. 5 pro 100.000 EW)
    \item[\gT{Letalit\"at}]  \index{Letalit\"at}Anteil der
    an XY Verstorbenen beozgen auf die Gesamtzahl der Erkrankten ({\quotedblbase}1 von 5 Erkranten
    stirbt{\textquotedblright})
\end{description}
\end{multicols}


\subsection{Entz\"undung}
\index{Entz\"undung}

\gDefinition{Reaktion des K\"orpers welche charakterisiert
ist durch:}

\begin{enumerate}
    \item \gEs{R\"otung}
    \item \gEs{Schwellung}
    \item \gEs{\"Uberw\"armung}
    \item \gEs{Schmerz}
    \item \gEs{Funktionsverlust}
\end{enumerate}

\paragraph{Ausl\"oser}

\begin{itemize}
    \item Fremdk\"orper
    \item Krankheitserreger
    \item Mechanische Beanspruchung
    \item ...
\end{itemize}

\paragraph{Anh\"angsel
{\textquotedblleft{}\gT{-itis}{\textquotedblright}}}\index{{}-itis}
Gastritis ${\rightarrow}$
{\quotedblbase}Magenentz\"undung{\textquotedblleft}, Hepatitis
${\rightarrow}$ Leberentz\"undung, Pankreatitis ${\rightarrow}$
Entz\"undung der Bauchspeicheldr\"use, Meningitis ${\rightarrow}$
Hirnhautentz\"undung, Rhinitis ${\rightarrow}$
{\quotedblbase}Schnupfen{\textquotedblleft}, Otitis ${\rightarrow}$
Ohrenentz\"undung, ...

\paragraph{Die wichtigsten
\"Ubertragungswege}
\begin{itemize}
    \item f\"akal-oral (Schmierinfektion, zB. Badesee)
    \item Aerogen (Tr\"opfcheninfektion)
    \item Genital (Geschlechtsverkehr)
    \item \"Uber/durch die Haut
    \item Diaplazent\"ar (w\"ahrend Schwangerschaft)
    \item W\"ahrend der Geburt
    \item Eintrittspforten: alles was irgendwie offen ist
    \begin{itemize}
        \item Wunden, Magen-Darm-Trakt, Atmung, Auge, Penis/Scheide, Plazenta
        \ldots
    \end{itemize}
    \item Fremdk\"orper sind Wohnh\"auser und Autobahnen f\"ur Keime
    \begin{itemize}
        \item Harnkatheter, ven\"ose Zug\"ange, Tubus, k\"unstlicher
        Mageneingang, Dialysekatheter {\dots}
    \end{itemize}
\end{itemize}
\paragraph[]{}
\paragraph[Kreuzinfektion]{Kreuzinfektion}\marginlabel{Kreuzinfektion}

\gAuthNotes{GABS: Definition!}
\index{Kreuzinfektion}Kreuzinfektion: \gEs{Patient 1
${\rightarrow}$ \emph{Sanit\"ater} ${\rightarrow}$ Patient
2}

\subsection{Die wichtigsten Erreger}

\begin{itemize}
    \item Bakterien
    \begin{itemize}
        \item klein, im Mikroskop meist sichtbar;
        {\textquotedblleft}Normale{\textquotedblright} Zellen
    \end{itemize}
    \item Viren
    \begin{itemize}
        \item sicher klein, im Mikroskop i.d.R. nicht sichtbar, keine echten
        Zellen, m\"ussen Zellen befallen
    \end{itemize}
    \item Pilze   
    \item Tiere, Tierparasiten
    \begin{itemize}
        \item W\"urmer 
    \end{itemize}
    \item (Toxine)
\end{itemize}


\subsubsection{Bakterien}

\index{Bakterien}

\begin{itemize}
    \item {\quotedblbase}Echte{\textquotedblleft} Zellen, die leben und
    aktiv etwas tun
    \begin{itemize}
        \item zB. Produktion von Substanzen (Toxine ...), Gasen, ... 
    \end{itemize}
    \item Viele Bakterien k\"onnen f\"ur den Menschen gut, egal oder auch
    b\"ose sein
    \item Haben die unterschiedlichsten Formen
    \begin{itemize}
        \item {\textquotedblleft}Kugeln{\textquotedblright}, St\"abchen, mit
        Gei{\ss}eln ...
    \end{itemize}
\end{itemize}


\paragraph{Bakterien und Antibiotika}

\begin{itemize}
    \item Antibiotika t\"oten prinzipiell Bakterien, aber:
    \begin{itemize}
        \item Antibiotika haben ein bestimmtes Spektrum, dh. wirken nur auf
        bestimmte Keime
        \item Bakterien entwickeln Resistenzen gegen Antibiotika
        \item Antibiotika helfen dann nicht mehr
    \end{itemize}
\end{itemize}

\gFazit{Resistenzen sind ein enorm gro{\ss}es Problem in der
heutigen Medizin!

\ldots und damit auch im Rettungs- und Krankentransport.}

{\textquotedblleft}\index{Spitalskeime}\gT{Spitalskeime}{\textquotedblright}:
Keime mit {\textquotedblleft}angez\"uchteten{\textquotedblright}
Resistenzen

Unterscheide:
{\textquotedblleft}Wald-und-Wiesen{\textquotedblright}-Keime und
Spitalskeime 

\gFazit{${\rightarrow}$ Keim ist nicht gleich Keim, auch
wenn sie den gleichen Namen haben}

\gMarried
{\paragraph{Bakterien und Pathogenit\"at}~

Die B\"osartigkeit von Bakterien h\"angt u.a. ab von:
    \begin{itemize}
        \item Welches Bakterium
        \item Spezies, Stamm, Resistenzen
        \item Wo ist die Infektion
        \item Wie besiedelt ist die Umgebung
        \item Wie ist der Immunstatus des Patienten
        \item Kind, Chemotherapie, AIDS, ... 
        \item K\"orperliche Verfassung
        \item Alter, Ern\"ahrung, Stress, ...
    \end{itemize}
}{~

\begin{center}
    \includegraphics[width=2.858cm]{images/AASdev20090228-img96.jpg}
    \captionof{figure}{F\"ur manche Bakterien zahlt man
    freiwillig viel Geld, hier am Beispiel des Lacotbacillus casei gezeigt. }
\end{center}}

Daraus folgt:
    \begin{itemize}
        \item Manche Keime machen immer krank
        \item Manche Keime nur unter Umst\"anden
        \item \gEs{Opportunisten} nur wenn der Wirt
        geschw\"acht ist, und das sind viele Patienten im Rettungs- und Krankentransport!
    \end{itemize}





\subsubsection{Viren}\marginlabel{Viren}
\index{Viren}\begin{itemize}
\item Keine Zellen, sehr klein, im Lichtmikroskop nicht sichtbar
\item Brauchen Wirtszelle um sich zu vermehren
\item Schleusen eigenes genetisches Programm in die Wirtszelle ein
\item zB: viele Verk\"uhlungen, Grippe, Feuchtblattern, viele
Hepatitis-Formen ...\item
\gZ{
\item Nachweis:
\begin{itemize}
    \item NICHT anz\"uchtbar auf N\"ahrb\"oden
    \item Antigen- oder Antik\"orpernachweis
    \item PCR (DNA/RNA-Nachweis)
\end{itemize}
}
\end{itemize}




\section{Wichtige
Krankheitsbilder}\marginlabel{Krankheitsbilder}


\subsection{Grippe -- Influenza}
\index{Grippe}\index{Influenza}\index{Influenza}

\gDefinition{Eine schwere, durch Viren verursachte
Infektionskrankheit.}

Die Grippe ist eine \gE{virale} Erkrankung, welche je nach
Patient auch einen sehr schweren Verlauf nehmen kann und
\gEs{Komplikationen} wie Lungen- oder
Herzmuskelentz\"undungen oder weitere Infektionen nach sich ziehen kann. Besonders gef\"ahrdet sind \"altere und immunsupprimierte
Personen. Gegen die Grippe wird eine \gEs{Schutzimpfung}
angeboten. Diese wirkt f\"ur ein Jahr und muss jede Saison widerholt werden. Die Impfung
bietet keinen hundertprozentigen Schutz, wird aber trotzdem f\"ur
Risikopatienten und Personen die einem erh\"oten Infektionsrisiko
ausgesetzt sind (darunter f\"allt auch Personal im Gesundheitsbereich)
empfohlen. 

\gFazit{Da es sich bei der Grippe um eine durch Viren
verursachte Krankheit handelt sind \gEs{Antibiotika
wirkungslos}.}

Die Grippe tritt besonders in der kalten Jahreszeit auf und tritt
regelm\"a{\ss}ig als \index{Epidemie}\gEs{Epidemie} auf. Es
kann auch zum Ausburch einer weltweiten \index{Pandemie}Pandemie kommen, \gZ{z.B.
Spanische Grippe (1918, 25~Mio~Tote), Asiatische Grippe
(1957, 1~Mio~Tote) und  Hongkong-Grippe (1968,
800.000~Tote)\footnote{http://de.wikipedia.org/wiki/Influenza}}.

In \"Osterreich erkranken pro Jahr ca. 380.000 Personen an der
Influenza, ca. 4.500 Personen m\"ussen station\"ar behandelt werden und
ca. 120 Personen sterben an den
Folgen\footnote{http://www.bmg.gv.at/cms/site/standard.html?channel=CH0742\&doc=CMS1075899626901}.

\gFazit{Die Grippe tritt in Wellen auf und kann t\"odlich
enden.}


\paragraph{Symptome}

Typisch f\"ur die Grippe ist ist das \gEs{rasche, hohe
anfiebern (38{\textdegree}-40{\textdegree}C) mit Sch\"uttelfrost,  Kopf-, Hals-,
Muskel-, Gelenksschmerzen, Husten und stark reduziertem
Allgemeinszustand} sowie langdauernde Abgeschlagenheit nach
Ende der Erkrankung. Bei schon vorher gschw\"achten Patienten kann der Verlauf
weniger spektakul\"ar aussehen (z.B. kaum Fieber), das
Risiko einer Komplikationen ist allerdings nat\"urlich h\"oher. 

\subsection{Grippaler Infekt}\index{grippaler
Infekt}

\emph{Syn.: \index{Verk\"uhlung}Verk\"uhlung,
 \index{Erk\"altung}Erk\"altung} 

Die Verk\"uhlung ist eine Infektion der oberen Luftwege mit Viren,
allerdings mit anderen als bei der {\quotedblbase}echten
Grippe{\textquotedblleft} (Influenza). Die Symptome \"ahneln denen der
Grippe, allerdings sind sie bei weitem nicht so stark ausgepr\"agt. Im
Vordergrund stehen Halsschmerzen, Husten, rinnende oder verstopfte Nase
und eine eher leicht erh\"ohte Temperatur. 

\subsection{Wundstarrkrampf --
Tetanus}\index{Tetanus}\index{Wundstarrkrampf}

\emph{\gZ{Clostridium tetani}}

%\includegraphics[width=5.984cm,height=3.539cm]{images/AASdev20090228-img97.gif}
\captionof{figure}{ \gAuthNotes{Tetanus Bild in gif: img97} }

\begin{itemize}
    \item Der Erreger kommt \"uberall in der Umgebung vor
    \item Produziert Toxin durch welches die Skelettmuskulatur tonisch
    krampft
    \gZ{
    \item Inkubation: 4-21 Tage
    \item Vorzeichen: M\"udigkeit, Inappetenz
    }\gEs{
    \item Klinik: Tonische Kr\"ampfe der quergestreiften Muskulatur, oft
    beginnend mit Kaumuskulatur
    \item Problem: Schmerzen, Ateml\"ahmung
    \item Letalit\"at ca. 50\%
    \item Impfung!
    }
\end{itemize}

\paragraph{Tetanus Impfschutz}\index{Impfschutz}

\begin{itemize}
    \item Auch bei kleinen Verletzungen nach dem Impfschutz fragen und
    \gEs{dokumentieren!}
    \item Patient soll unverz\"uglich im Impfpass nachschauen.
    \item Wenn die letzte Impfung l\"anger als \gEs{5~Jahre}
    zur\"uck liegt oder der Patient noch nie geimpft wurde muss der Patient unverz\"uglich eine
    Unfallabteilung aufsuchen.
    \item Am Revers best\"atigen lassen!
\end{itemize}

\gCave{Eine mangelhafte Aufkl\"arung oder Dokumentation
bez\"uglich Tetanus und dessen Impfschutz ist grob fahrl\"assig und ein unverzeihbarer
Kunstfehler.}


\subsection{Tuberkulose}\marginlabel{TB}

\gZ{Mykobacterium Tuberculosis,}
"`\gT{Tbc}"'

\marginpar{\includegraphics[width=4cm]{images/AASdev20090228-img98.jpg}\captionof{figure}{Röntgenbild
eines TB-Pat\-ient\-en.
\gAuthNotesPicLegalNotOk{img98}{Durch Foto von Hauer
ersetzen}}}

\begin{itemize}
    \item Kann praktisch alle Organe befallen
    \item Der Verlauf wird vom Organbefall bestimmt.
    \begin{itemize}
        \item Lunge, Skelett, Darm, Haut, Genitalien, ...
    \end{itemize}
    \item Ausgangspunkt ist fast immer die \gEs{Lunge}.
    \item \"Ubertragung: aerogen, oral, Hautkontakt, plazent\"ar
    
    \item Inkubationszeit: 4-12 Wochen
    
    \item Klinik:
    \begin{itemize}
        \item \gEs{Subfebrile Temperatur, Nachtschwei{\ss},
        Husten, evt. Bluthusten}, Lymphknotenschwellung, 
        Pleuritis
        \item {\textquotedblleft}\index{Offene Tb}\gT{Offene
        Tb}{\textquotedblright}: Infektionsherd hat
        {\textquotedblleft}Luftkontakt{\textquotedblright} ${\rightarrow}$ Tr\"opfcheninfektion m\"oglich ${\rightarrow}$
        Schutzmaske
    \end{itemize}
    \item \gEs{Meldung an Leitstelle \& Betriebsleitung!}    
\end{itemize}



\subsection{HIV - Humanes Immundefizienz Virus}
\marginlabel{HIV}
\index{HIV}\index{Humanes Immundefizienz Virus}

Das HIV bef\"allt vor allem die Zellen des Abwehrsystems.
Es vermehrt sich in ihnen, setzt sie au{\ss}er Funktion und zerst\"ort sie schlie{\ss}lich. Das k\"orpereigene Abwehrsystem kann HIV nicht aus dem K\"orper
    entfernen.
Es kommt zu einer Immunschw\"ache.

HIV verursacht das AIDS (Aquired Immune Deficiency Syndrom,
erworbenes Immunschwäche-Syndrom). Dabei kommt es zu einer
Reihe von Begleiterkrankungen:

\begin{itemize}
    \item Infektionen mit normalerweise harmlosen Erregern
\begin{itemize}
    \item Pilze, {\textquotedblleft}fakultative{\textquotedblright} Erreger,
    {\textquotedblleft}Opportunisten{\textquotedblright}, seltene Arten der
    Lungenentz\"undung
\end{itemize}

    \item Tumore
\end{itemize}


\paragraph{Übertragung}

\begin{itemize}
    \item K\"orperfl\"ussigkeiten
    \item Blut, Blutprodukte, Sperma, Scheidensekret
    \item Ungesch\"utzter Sex, Transfusionen
    \item Nadelstichverletzungen
    \item KEINE \"Ubertragung durch Sozialkontakte (H\"andesch\"utteln,
    Haushalt, Geschirr, Toilettenanlagen)
\end{itemize}
\paragraph{Test}
\begin{itemize}
    \item Erst nach 6-12 Wochen sind Antik\"orper nachweisbar
    \item \gEs{HIV-Test nicht
    {\textquotedblleft}aktuell{\textquotedblright},
    {\textquotedblleft}12 Wochen blind{\textquotedblright}}
\end{itemize}
\paragraph{AIDS}
Nach ca 5-10 Jahren bricht AIDS aus   

\paragraph{Therapie}

\begin{itemize}
    \item Keine Heilung m\"oglich, aber deutliche Lebensverl\"angerung
\end{itemize}

\paragraph{Prophylaxe}

\begin{itemize}
    \item Postexpositionsprohylaxe bei Nadelstichverletzungen
    \item Kondome
\end{itemize}



\subsection{(Virale) Hepatitis}\index{Hepatitis} 

\gDefinition{Leberentz\"undung durch verschiedene (virale)
Erreger (A -- E)}

\paragraph{Impfung}Impfung nur f\"ur Hepatitis A und B
verf\"ugbar (Twinrix{\texttrademark})

Achtung: Non-Responder m\"oglich!

\paragraph{Klinik und Verlauf}

\begin{itemize}
    \item Vergr\"o{\ss}erte, schmerzhafte Leber
    \item Ikterus (Gelbsucht)
    \item Dunkler Urin (heller Sch\"uttelschaum),
    \item evt. heller Stuhl (Acholie)
    \item Unterscheide {\textquotedblleft}akut{\textquotedblright} und
    {\textquotedblleft}chronisch{\textquotedblright}
    \begin{itemize}
        \item Chronisch: wenn Hepatitis nicht abheilt
    \end{itemize}
    \item Ende: Leberversagen, Tod
\end{itemize}

\subsubsection{Hepatitis-A - Hepatitis-A-Virus (HAV)}

\begin{itemize}
\item \gZ{Inkubationszeit 10 - 40 Tage}  
\item typische Reisekrankheit (Mittelmeerraum, S\"udamerika, Orient),
meist epidemisch, 
\item Dauer einige Wochen, 
\item kein \"Ubergang in chronischen Verlauf
\end{itemize}


\subsubsection{Hepatitis-B - Hepatitis-B-Virus (HBV)}

\begin{itemize}
\item \gZ{Inkubationszeit 4 - 12 Wochen,} 
\item meist sporadisches Auftreten, 
\item \"Ubertragung durch Speichel, Blut(produkte), Sperma und andere
K\"orperfl\"ussigkeiten, 
\item chronischer Verlauf in ca. 5 \%
\end{itemize}


\subsubsection{Hepatitis C - Hepatitis-C-Virus (HCV)}

\begin{itemize}
\item \gZ{Inkubationszeit 6 - 12 Wochen,} 
\item h\"aufig nach Bluttransfusion oder i.v. Drogenabusus, 
\item \gEs{oft chronischer Verlauf} (in 70 -- 80~\% der
F\"alle)
\end{itemize}

\gFazit{F\"ur den Rettungsdienst sind besonders Hep B und C
wichtig: Hohe Infektiosit\"at, chronischer Verlauf!

Gegen Hepatitis C gibt es derzeit keine Impfung!}



\subsection{Sepsis}
\label{topic:sepsis}

Siehe auch: Septischer Schock und Toxischer Schock,
Kap.\,\ref{topic:schock-septisch}.

\begin{itemize}
\item Wenn Keime und deren Produkt ausser Kontrolle geraten erfolgt
Amoklauf des K\"orpers:

\item Ungehemmte Freisetzung von Stoffen des Entz\"undungs-, Gerinnungs-
und Immunsystems

\item Auch ohne Infektion m\"oglich:


\begin{itemize}
\item \gZ{Systemic Inflammatory Response Syndrome (SIRS)} 
\item nach schwerem Trauma, Verbrennungen, Gewebshypoxie
\end{itemize}\end{itemize}


\paragraph{Klinik} Klinik abh\"angig vom Fortschritt und
k\"orperlichen Abbau.


\begin{itemize}
\item Sehr stark reduzierter Allgemeinzustand

\item Fieber ({\textgreater}\,38{\textdegree} C) oder
Hypothermie ({\textless}\,36{\textdegree} C)

\item Herzfrequenz {\textgreater}\,90/min

\item Atemfrequenz {\textgreater}\,20/min oder
Hyperventilation

\item unkontrollierbare Gerinnung/Blutung

\item Hypotonie ${\rightarrow}$ Septischer Schock

\item Organversagen  

\end{itemize}

\subsection{Multiresistente Keime}

Vor der Entdeckung der Antibiotika waren bakterielle Infektionen sehr
gef\"urchtet, viele Menschen sind an -- nach heutiger Sicht simplen --
Infektionen gestorben. 

Heute kann man zwar solche Infektionen mit Antibiotika behandeln, jedoch
k\"onnen Bakterien \gEs{Resistenzen} gegen Antibiotika
entwickeln. Die sonst wirksamen Medikamente wirken dann nicht mehr oder nur noch sehr
abgeschw\"acht. Als behandelnder Arzt steht man dann hilflos vor dem
Patienten. 

\gFazit{Resistenzen machen Antibiotika wirkungslos.}

Es muss daher das Ziel aller sein die Ausbreitung von resistenten Keimen
zu verhindern. Dem Rettungs- und Krankentransport kommt dabei eine
\gE{Schl\"usselrolle} zu, da es hier leicht zu
\gEs{Kreuzinfektionen} kommen kann.


\gFazit{Resistenzen sind ein enorm gro{\ss}es Problem in der
heutigen Medizin!

\ldots und damit auch im Rettungs- und Krankentransport, da es
hier leicht zu Kreuzinfektionen kommen kann.}

{\textquotedblleft}\index{Spitalskeime}\gT{Spitalskeime}{\textquotedblright}:
Keime mit {\textquotedblleft}angez\"uchteten{\textquotedblright}
Resistenzen

\gFazit{Unterscheide:
{\textquotedblleft}Wald-und-Wiesen{\textquotedblright}-Keime und
Spitalskeime} 

Der wohl bekannteste Problemkeim ist der \gT{MRSA}
(\gZ{Multiresistenter Staphylococcus aureus, siehe unten)}.
Das liegt schlicht daran dass er der erste war. In den letzten Jahren ist die Resistenzentwicklung aber
auch bei vielen anderen Keimen zu beobachten, \gZ{Beispiele
daf\"ur sind die ESBL{\textquotesingle}s und VRE{\textquotesingle}s sowie der VRSA
(Extended Spectrum Beta-Laktamase-Bildner, Vancomycin-resistente
Enterokokken, Vancomycin-resistenter Staph. aureus)}.
\gEs{Das Problem der multiresistenten Keime wird immer
gr\"o{\ss}er.}

\gFazit{Je wirkungsloser die Antibiotika werden desto
wichtiger ist die Hygiene und die Vermeidung von
(Kreuz-)Infektionen.}

\begin{minipage}{13.67cm}
[Warning: Draw object ignored]

\captionof{figure}{Teufelskreis der
resistenten Keime. Oft das Bindeglied: Der Rettungs- und
Krankentransportdienst.
\gSourcePicture{unbekannt} \gAuthNotes{Achtung! Zeichnung
fehlt!}}


\end{minipage}

\subsubsection{MRSA -- Multi-Resistenter Staphylococcus
aureus} \gZ{Urspr\"unglich: Methicilin-Resistenter Staphylococcus
Aureus}

\gMarried
{Der \emph{Staphylococcus aureus} ist ein
"`Allerweltskeim", bis zu 50\% der
Bev\"olkerung sind besiedelt. Vor allem ist die vordere
Nasenh\"ole besiedelt, aber auch die Haut (Achselh\"ole,
Analregion) und die Schleimhäute können ihn beherbergen.
Auch im Spital ist er häufig, aufgrund des hohen
Antibiotikaverbauchs ist er dort aber oft schon
Antibiotika-Resistent: Man spricht dann vom \gT{MRSA}.}
{~

\begin{itemize}
    \item "`Allerweltskeim"'
    \item aber resistente Form:
    \gT{MRSA}\footnote{Multi-resistenter Staphylococcus aureus}
    \begin{itemize}
        \item {\textquotedblleft}Spitalsinfektionen{\textquotedblright}
        \item schwer bis gar nicht behandelbar
        \item Wunden, die nicht abheilen wollen, Sepsis, Tod ...
    \end{itemize}
\end{itemize}
}






\subparagraph{MRSA-Transport -- Allgemeine Ma{\ss}nahmen}
\label{bkm:RefHeading40262819}

\begin{itemize}
    \item Informationen \"uber MRSA-Tr\"agerschaft
    \marginlabel{gilt im Prinzip auch f\"ur die anderen
    {\quotedblbase}Problemkeime{\textquotedblleft}}
    \item Art/Ort der MRSA-Besiedlung
    \item Patientenvorbereitung und Transport
    \item Wunden sind frisch verbunden und gut abgedeckt
    \item Bei Besiedlung der Atemwege tr\"agt der Patient einen Mund-
    /Nasenschutz
    \item Vor dem Transport f\"uhrt der Patient eine hygienische
    H\"andedesinfektion durch
    \item Allgemeine Hygienema{\ss}nahmen
    \begin{itemize}
        \item Einmalhandschuhe und Schutzkittel.
        \item Nach Ablegen sofortige hygienische H\"andedesinfektion
        \begin{itemize}
            \item Nach Abschluss des Patiententransportes sind alle Materialien,
            Ger\"ate, Instrumente und Fl\"achen, welche direkten Kontakt mit dem
            Patienten hatten gem\"a{\ss} dem \"ublichen Hygieneplan zu
            desinfizieren.
        \end{itemize}
    \end{itemize}
    \item Kugelschreiber!!!
    \item Alle waagerechten Oberfl\"achen des Fahrzeuginnenraumes sind mit
    einem Fl\"achendesinfektionsmittel in Form einer
    Scheuerwischdesinfektion desinfizierend zu reinigen.
    \item Abschlie{\ss}end f\"uhrt das Einsatzpersonal eine hygienische
    H\"andedesinfektion durch.
    \item Nach Abschluss der vorangegangenen Hygiene- und
    Desinfektionsma{\ss}nahmen ist das Fahrzeug wieder voll einsatzbereit.
\end{itemize}

F\"ur die anderen Problemkeime \gZ{(VRSA, VRE, ESBL,
{\dots})} gelten sinngem\"a{\ss} die gleichen
Verfahrensweisen.





\section[Desinfektion]{Desinfektion}
\index{Desinfektion}\begin{center}
                    
                    
                    \end{center}
\begin{tabularx}{\textwidth}[H]{XX}
\toprule
\gT{Desinfektion}
 &
\gT{Sterilisation}
\\\midrule


Ziel: \gEs{Keim\underline{reduktion}} ${\rightarrow}$ nicht
infekti\"os

Richtet sich gegen ausgew\"ahlte Keime

Aber: nicht keimfrei! nicht steril!

Chemische Substanzen

Einlegebad

 &

Abt\"otung aller Mikroorganismen -
{\textquotedblleft}\gEs{\underline{keimfrei}}{\textquotedblright}

Dampfsterilisation (\"Uberdruck)

Hei{\ss}luftsterilisation

\\\bottomrule
\end{tabularx}


\paragraph{Hygieneplan -
Hygienebeauftragter}\index{Hygieneplan}~

\gMarried{Der Hygieneplan regelt die konkreten Hygienema{\ss}nahmen und Was Wann
Wie Womit (Wo) zu geschehen hat, siehe allgemeiner Aushang.}
{Der derzeitige Hygienebeauftragte ist

~

\largepencil \hrulefill}.

\subsubsection{Pers\"onliche Hygiene}
\marginpar{\includegraphics[width=4cm]{images/AASdev20090228-img99.jpg}}

\begin{itemize}
    \item Schmuck, Uhren, Freundschaftsarmb\"ander sind perfekte
    N\"ahrb\"oden ${\rightarrow}$ Vor dem Dienst ablegen!
    \item Dienstkleidung regelm\"a{\ss}ig waschen!
    \item Bei Kontamination wechseln!
    \item Verwendung von Plastiksch\"urzen
    \item Bei Dienstantritt
    \begin{itemize}
        \item H\"ande waschen
        \item Nagelfalze reinigen
        \item Hygienische H\"andedesinfektion
        \index{H\"andedesinfektion}\index{Hygienische H\"andedesinfektion}
    \end{itemize}
\end{itemize}



\subsubsection{Flächen- und Wischdesinfektion}

\begin{itemize}
    \item Mit Fl\"achendesinfektionsmittel \index{Fl\"achendesinfektionsmittel}
    \item Einwirkzeit lt. Packung beachten!
    \item Fl\"ache muss w\"ahrend der Zeit feucht sein
    \item \gEs{in Autos nur Fl\"achendesinfektionsmittel
    ohne Alkohol!}
    \begin{itemize}
        \item Alkohol zerst\"ort Acryl-Fl\"achen! ${\rightarrow}$
        \index{\gT{Acryl-Des}{\texttrademark}}Acryl-Des{\texttrademark}
    \end{itemize}
    \item nicht den Schmutz \"uberall verteilen
    \begin{itemize}
        \item \gEs{Von au{\ss}en in Richtung Schmutzquelle
        wischen}
        \item Einmalt\"ucher wechseln
    \end{itemize}
\end{itemize}

\gCave{In den Fahrzeugen d\"urfen keine alkoholischen
Desinfektionsmittel verwendet werden da diese die
Acryloberfl\"achen zerst\"oren k\"onnen!}

\subsubsection{Ger\"atedesinfektion}

\begin{wrapfigure}{r}{4cm}
\includegraphics[width=4cm]{./images/AASdev20090228-img100.jpg}
    \captionof{figure}{Desinfektionswannen, verschiedene
    Gr\"o{\ss}en}
\end{wrapfigure}

\begin{itemize}
    \item Nach Kontamination:
    \begin{itemize}
        \item \gEs{Grobe Reinigung vor Ort}: Eintrocknen
        verhindern
        \item Wischdesinfektion
        \item Ggfs. Schl\"auche durchsp\"ulen
        \item Fl\"achendesinfektionsmittel in Nierenschale / Joghurtbecher,
        saugen ...
    \end{itemize}
    \item Meldung an Leitstelle:
    \emph{{\textquotedblleft}Nicht EB, m\"ussen
    putzen{\textquotedblright}}
    \item Kontaminierte Teile gegen Ersatzmaterial tauschen
    \item Zerlegen
    \item Nochmals sp\"ulen
    \item Wenn Desinfektionswanne leer ${\rightarrow}$ einlegen
    \begin{itemize}
        \item \gEs{Vollst\"andig zerlegt} laut Anleitung
        \item \gEs{Luftblasenfrei!}
        \item Mit Niederhaltesieb unter die Oberfl\"ache dr\"ucken
        \item Zeit auf Zettel notieren
        \begin{itemize}
            \item {\textquotedblleft}was, wann, von wo, von wem
            eingelegt{\textquotedblright}
        \end{itemize}
        \item sonst nur bereit legen, in Plastiksack
        \begin{itemize}
            \item {\textquotedblleft}was, von wo, von wem, warum{\textquotedblright}
        \end{itemize}
    \end{itemize}
\end{itemize}


\subsubsection{Hygienische H\"andedesinfektion}
Ziel:
\begin{itemize}
    \item Keimreduktion
    \item Vermeidung von \gEs{Kreuzinfektionen}
    \begin{itemize}
        \item Patient1 ${\rightarrow}$ Sanit\"ater ${\rightarrow}$ Patient 2
        \item \gEs{{\textquotedblleft}Unterbrechung der
        Infektionskette{\textquotedblright}}
    \end{itemize}
\end{itemize}
Durchzuf\"uhren:
\begin{itemize}
    \item Dienstantritt
    \item \gEs{Nach jedem Patienten!!}
\end{itemize}
Womit?
\begin{itemize}
    \item H\"andedesinfektionsl\"osung: Sterilium, Monopronto
\end{itemize}

\begin{wrapfigure}{r}{5cm}
\includegraphics[width=4.001cm]{images/AASdev20090228-img101.jpg}
\end{wrapfigure}

\minisec{Vorgehen} 

\begin{enumerate}
    \item Ringe entfernen
    \item Eine Portion alkoholisches H\"andedesinfektionsmittel (ca. 3ml =
    1~Hohlhand) aus dem Spen\-der entnehmen.
    \item mittels Ellenbogen, vollkommene Benetzung beider H\"ande und
    Handgelenke
    \item Handinnenfl\"achen gegeneinander reiben
    \item geschlossene Finger
    \item Handgelenke einreiben
    \item Mit rechter Handinnenfl\"ache linken Handr\"ucken reiben und
    umgekehrt, dabei Finger ineinander verschr\"anken.
    \item Mit ineinander verschr\"ankten Fingern Handinnenfl\"achen
    gegeneinander reiben.
    \item H\"ande ineinander verhaken und Finger gegeneinander bewegen
    \item Daumen mit gegen\"uberliegender Hand vollst\"andig umschlie{\ss}en
    und rotierend reiben.
    \item Daumenkuppen nicht vergessen!
    \item Fingerkuppen im Handteller kreisf\"ormig einreiben bis Alkohol
    verdunstet ist
    
\end{enumerate}
 Einwirkzeit von mindestens 30 Sekunden beachten! 
 
\begin{figure}[H]
\includegraphics[width=3.961cm,height=3.169cm]{images/AASdev20090228-img102.jpg}
\end{figure}


\subsubsection{\"Offnen von sterilem Material}

Packung {\quotedblbase}aufsch\"alen{\textquotedblleft} 

\begin{figure}[H]
\includegraphics[width=6.656cm,height=4.993cm]{images/AASdev20090228-img103.jpg}
\includegraphics[width=6.663cm,height=4.984cm]{images/AASdev20090228-img104.jpg}
\caption{Öfnnen von sterilen Material unter verschiedenen
Bedingungen: Immer wird die Verpackung "`aufgeschält"'.
\gAuthNotesPicLegalNotOk{}{}}
\end{figure}

\subsubsection{Selbstschutz}

\begin{itemize}
    \item Immer Einmalhandschuhe beim Patientenkontakt tragen!
    \item Handschuhe nach jedem Patienten wechseln
    \item Nie mit schmutzigen Handschuhen in die Tasche greifen!
    \item \gEs{Ein Helfer hat Patientenkontakt, der andere
    reicht zu!}
    \item Achtung {\textquotedblleft}Kugelschreiber{\textquotedblright}
    \item Einmalsch\"urzen bei Bedarf
    \begin{itemize}
        \item evt. Infektionsschutzanz\"uge, Maske
    \end{itemize}
    \item \gEs{Spitze Gegenst\"ande (Nadeln, Lanzetten, ...
    sofort in die Kontamed-Box!}
    \begin{itemize}
        \item Ausnahme: ausdr\"uckliche Arztanweisung (Blutzuckerbestimmung aus
        Nadel)
        \item \emph{{\textquotedblleft}Wer sticht, der
        entsorgt!{\textquotedblright}}
    \end{itemize}
\end{itemize}

\paragraph{Kontamed, Sharp, Gelbe Box ...}
Spitze Gegenst\"ande (Nadeln, Lanzetten, ... sofort in Gelbe Box!

Ausnahme: ausdr\"uckliche Arztanweisung (Blutzuckerbestimmung aus Nadel)

\gFazit{Wer sticht, der entsorgt!}



\begin{figure}[H]
\includegraphics[width=3.377cm,height=4.552cm]{images/AASdev20090228-img105.jpg}
\includegraphics[width=5.71cm,height=4.289cm]{images/AASdev20090228-img106.jpg}
\caption{Kontamed-Boxen, verschiedene Fabrikate.}
\end{figure}

\gFazit{NIE Nachstopfen.

NIE ausleeren.

Wenn die Box voll ist wird sie verschlossen und abgegeben.}

\subsubsection{Wundreinigung}

\begin{itemize}
    \item Bei frischen, nicht infizierten Wunden:
    \begin{itemize}
        \item Reinigung von innen nach au{\ss}en
        \item Ziel: keine Keime in Wunde bringen
    \end{itemize}
    \item Womit?
    \begin{itemize}
        \item Sterilen L\"osungen oder Octenisept
        \item Sterile Tupfer
    \end{itemize}
\end{itemize}

KEINE Papierhandt\"ucher oder unsterile Tupfer !!!

\begin{itemize}
    \item Danach steriler Verband!
    \item Bei offenen Frakturen ist keimfreiheit be\-sonders wichtig !!!
    \begin{itemize}
        \item Gefahr einer lebenslangen Knochen\-marks\-ent\-z\"undung
        (Osteomyelitis)!
        \item \index{OpSite}\gT{OpSite}\texttrademark-Folie
        wenn vorhanden.
    \end{itemize}
\end{itemize}

\begin{figure}[H]
    \includegraphics[width=5.742cm,height=5.658cm]{images/AASdev20090228-img107.png}
    \caption{Reinigung einer
    sauberen Wunde: Von Innen nach Aussen, der Tupfer wird anschliessend
    verworfen. \gAuthNotesPicLegalNotOk{}{}}
\end{figure}
\begin{figure}[H]
    \includegraphics[width=4.219cm,height=4.219cm]{images/AASdev20090228-img108.jpg}
    \caption{Anbringen
    einer OpSite(TM)-Folie: Es ist wichtig die Folie zu spannen damit sie
    sich nicht selbst verklebt.
    \gAuthNotesPicLegalNotOk{}{} }
\end{figure}



\section{Besondere Verhaltensweisen}
\subsection{Nadelstichverletzungen}\index{Nadelstichverletzungen}
\begin{itemize}
    \item Erstma{\ss}nahmen durchf\"uhren:
    \begin{itemize}
        \item bei Stich-bzw.Schnittverletzungen: Wunde auspressen und 10 Minuten
        desinfizieren
        \item bei Spritzern mit Blut oder Sekreten in Auge oder Mund:
        Schleimhaut mit Wasser u./o. Schleimhautdesinfektionsmittel sp\"ulen
    \end{itemize}
    \item Leitstelle informieren
    \item ${\rightarrow}$ \gEs{Umgehend in ein Unfallspital}\footnote{am besten der der AUVA, da Arbeitsunfall}
    \item Patient: Zielspital informieren:
    \begin{itemize}
        \item Blut sichern f\"ur HIV- und Hepatitis-Serologie
    \end{itemize}
    \item  \gC{Unfallmeldung}  (wichtig!), Blutabnahme
    \item ggf. HIV-Postexpositionsprohylaxe (PEP)
    \begin{itemize}
        \item kurzfristige Therapie (4 Wo)
        \item Einzelfall muss immer individuell beurteilt werden
        (Nutzen/Risiko), starke Nebenwirkungen.
        \item Muss innerhalb von 2 Std. begonnen werden ${\rightarrow}$ daher
        SOFORT in ein Unfallspital der AUVA
    \end{itemize}
    \item Meldung an Betriebsleitung
    \item \gEs{Unfallmeldung (wichtig!) an die AUVA}
\end{itemize}

\gCave{Unfallmeldung nicht vergessen!

Sofort in ein Unfallspital! }


\subsubsection{MRSA und andere Resistente Keime}
\index{MRSA}

Die Ma{\ss}nahmen beim Transport von Patienten mit resistenten Keimen
werden unter {\quotedblbase}MRSA-Transport - Allgemeine
Ma{\ss}nahmen{\quotedblbase} 
\ref{bkm:RefHeading40262819} auf
Seite \pageref{bkm:RefHeading40262819} besprochen.
